\documentclass[reprint,aps,prd,letterpaper,superscriptaddress,nofootinbib,longbibliography,floatfix]{revtex4-2}

%% ----------------------------------------------------------------
%%  PACKAGES
%% ----------------------------------------------------------------
\usepackage[utf8]{inputenc}
\usepackage[T1]{fontenc}
\usepackage{amsmath,amssymb,amsthm}
\usepackage{mathtools}
\usepackage{bm}
\usepackage{booktabs}
\usepackage{xcolor}
\usepackage{graphicx}
\usepackage{microtype}
\usepackage{hyperref}

\hypersetup{
  hidelinks,
  pdftitle={Planck-Boundary Geometry of Compact S3 Cosmology},
  pdfauthor={Boris Batenin and Andrei Preece}
}

\emergencystretch=2em
\allowdisplaybreaks

%% ----------------------------------------------------------------
%%  THEOREM ENVIRONMENTS  (with epistemic labelling)
%% ----------------------------------------------------------------
\newtheorem{theorem}{Theorem}[section]
\newtheorem{corollary}[theorem]{Corollary}
\newtheorem{proposition}[theorem]{Proposition}
\newtheorem{lemma}[theorem]{Lemma}
\newtheorem{definition}[theorem]{Definition}

\theoremstyle{remark}
\newtheorem{remark}[theorem]{Remark}
\newtheorem{assumption}[theorem]{Assumption}
\newtheorem{ansatz}[theorem]{Ansatz}
\newtheorem{interpretation}[theorem]{Interpretation}

%% ----------------------------------------------------------------
%%  MACROS
%% ----------------------------------------------------------------
\newcommand{\lP}{\ell_P}
\newcommand{\tP}{t_P}
\newcommand{\mP}{m_P}
\newcommand{\zetaS}{\zeta_{S^3}}
\newcommand{\Riemann}{\zeta_{\rm R}}
\newcommand{\Czeta}{C_\zeta}
\newcommand{\model}{\texorpdfstring{\(\mathrm{FB}(S^3)\mathrm{R}\)}{FB(S3)R}}
\newcommand{\modelmath}{\mathrm{FB}(S^3)\mathrm{R}}

%% ----------------------------------------------------------------
\begin{document}
%% ----------------------------------------------------------------

\title{Planck-Boundary Geometry of Compact \texorpdfstring{$S^3$}{S3} Cosmology}

\author{Boris Batenin}
\affiliation{Independent Researcher, Berlin, Germany}

\author{Andrei Preece}
\affiliation{Independent Researcher, Berlin, Germany}

\date{May 2026}

\begin{abstract}
We analyze a compact \(S^3\) cosmological framework in the limiting regime
\(a\to\ell_P\).  The Planck epoch is modelled as a regular self-adjoint
boundary of minisuperspace, not as a classical curvature singularity: the
effective domain is \(a\in[\ell_P,\infty)\), the wave function may have a finite
non-zero value at \(a=\ell_P\), and the sub-Planck interval is not assigned a
probability measure in the effective theory.

The model is restricted to the smooth closed connected orientable
simply-connected spatial sector.  In that sector, the Poincar\'e--Perelman
theorem, together with uniqueness of smooth structures on three-manifolds,
selects \(\Sigma\simeq S^3\).  The scalar spectrum of the round three-sphere is
then organized by a Fibonacci coarse-graining.  The golden ratio
\(\varphi=(1+\sqrt5)/2\) enters only as the asymptotic ratio of successive
Fibonacci scales; it is not an additional dynamical constant.

Within the stated assumptions we prove four conditional results.  First, for a
closed FLRW metric with \(k=+1\), the Kretschmann scalar satisfies
\(K(t_P)\le C/\ell_P^4<\infty\) whenever the regularised scale factor is bounded
below and \(H,\dot H\) are finite at the Planck boundary.  This is a kinematic
semiclassical regularity theorem, not a microscopic quantum-gravity derivation.
Second, for any scale-covariant compact \(S^3\) spectrum whose spectral zeta
function is finite at the Casimir point,
\[
  \rho_{\rm vac}^{\rm reg}(a)=\frac{\Czeta}{2\pi^2a^4}.
\]
For the conformally shifted scalar reference spectrum,
\(\Czeta=1/240\), so
\(\rho_{\rm vac}^{\rm reg}=1/(480\pi^2a^4)\).  This is a regularised
compact-spectrum contribution, not a prediction of the observed late-time dark
energy density.  Third, a self-adjoint Wheeler--DeWitt operator on
\(L^2([\ell_P,\infty),a\,da)\) gives zero probability flux through the Planck
boundary.  Fourth, diffeomorphism invariance gives Noether-II differential
identities; in covariant phase space the full closed Cauchy slice has no
Iyer--Wald boundary Hamiltonian variation because \(\partial\Sigma=\varnothing\),
although quasi-local charges on internal subdomains remain possible.
\end{abstract}

\keywords{Planck boundary; closed cosmology; \(S^3\) topology; Fibonacci spectral coarse-graining; spectral zeta regularisation; Wheeler--DeWitt equation; Noether's second theorem; bounded curvature; compact Casimir energy}

\maketitle

%% ================================================================
\section{Introduction and Scope}
%% ================================================================

Classical FLRW cosmology, extrapolated without modification to \(a=0\), produces
divergent curvature invariants.  Perturbative quantum field theory, extrapolated
without regularisation to arbitrarily short wavelengths, produces ultraviolet
divergences.  The Planck length
\[
  \ell_P=\sqrt{\hbar G/c^3}
\]
marks the scale at which neither extrapolation should be taken as literal
classical physics.

This paper isolates a conservative Planck-boundary structure.  It does not claim
to construct a full microscopic theory of quantum gravity.  The framework has
five ingredients:
\begin{enumerate}
\item a spatial sector assumption: smooth, closed, connected, orientable, and
  simply connected three-manifolds;
\item a closed FLRW homogeneous geometry with \(k=+1\);
\item a lower effective boundary \(a=\ell_P\) for minisuperspace;
\item a compact \(S^3\) spectral zeta regularisation of vacuum sums;
\item a self-adjoint Wheeler--DeWitt boundary condition at \(a=\ell_P\).
\end{enumerate}
The mathematical question is then precise: what follows from these assumptions,
and what does not?

The answer is deliberately limited.  The model gives a conditional curvature
bound, a finite regularised compact-spectrum scale, norm preservation in the
minisuperspace sector, and the expected Noether-II identities of generally
covariant dynamics.  It does not derive finite \(H(t_P)\) or \(\dot H(t_P)\)
from first principles, does not predict the observed cosmological constant, and
does not turn the Fibonacci hierarchy into an exact spacetime symmetry.

\medskip
\noindent\textbf{Conventions.}
We use natural units \(c=\hbar=k_B=1\) except when defining \(\ell_P\).
The reduced Planck mass is \(\mP=(8\pi G)^{-1/2}\), the metric signature is
\((-+++) \), and \(S^3\) denotes the unit three-sphere.  The physical radius is
\(a(t)\).

%% ================================================================
\section{Topological and Spectral Framework}
%% ================================================================

\subsection{Selection of the \texorpdfstring{$S^3$}{S3} sector}

The topology is not derived from normalisability of the wave function.  It is a
sector assumption.

\begin{assumption}[Closed simply-connected spatial sector]
\label{ass:closed-sector}
The admissible spatial slice is restricted to the smooth, closed, connected,
orientable, simply-connected three-manifold sector:
\[
  \Sigma\ \text{smooth, closed, connected, orientable},
  \qquad \pi_1(\Sigma)=0 .
\]
\end{assumption}

\begin{theorem}[Topological uniqueness in the assumed sector]
\label{thm:topology}
Under Assumption~\ref{ass:closed-sector}, the spatial slice is homeomorphic to
\(S^3\).  In the smooth three-dimensional category used here, it is
diffeomorphic to the standard \(S^3\):
\[
  \Sigma \simeq S^3 .
\]
\end{theorem}

\begin{proof}
The Poincar\'e--Perelman theorem states that every closed simply-connected
topological three-manifold is homeomorphic to \(S^3\)
\cite{Perelman1,Perelman2,Perelman3}, building on Hamilton's Ricci-flow
programme~\cite{Hamilton}.  Moise's theorem gives the uniqueness, up to
diffeomorphism, of smooth structures on topological three-manifolds
\cite{Moise}.  Hence in the smooth category relevant for the metric formulation
the assumed sector is the standard \(S^3\) sector.
\end{proof}

\begin{remark}
Compactness supports a global \(L^2\) inner product and orientability supplies a
global volume form.  Simple connectivity excludes lens spaces and other closed
topological sectors.  These are assumptions of the present model; they are not
deduced from the Wheeler--DeWitt equation alone.
\end{remark}

\subsection{Fibonacci coarse-graining of the \texorpdfstring{$S^3$}{S3} spectrum}

The scalar Laplacian on the unit \(S^3\) has eigenvalues
\[
  -\Delta_{S^3}Y_{\ell mn}=\ell(\ell+2)Y_{\ell mn},
  \qquad \ell=0,1,2,\ldots ,
\]
with eigenspace dimension \((\ell+1)^2\).

Let
\[
  F_0=0,\qquad F_1=1,\qquad F_{n+1}=F_n+F_{n-1}\quad(n\ge1)
\]
be the Fibonacci sequence.  To avoid the repeated cut point \(F_1=F_2=1\), define
a strictly increasing cut sequence by
\[
  \mathcal F_0=0,\qquad
  \mathcal F_1=1,\qquad
  \mathcal F_n=F_{n+1}\quad(n\ge2).
\]
The non-empty Fibonacci spectral bands are
\[
  B_n=\{\ell\in\mathbb N_0:\mathcal F_n\leq \ell<\mathcal F_{n+1}\},
  \qquad n=0,1,2,\ldots .
\]
Thus \(B_0=\{0\}\), \(B_1=\{1\}\), \(B_2=\{2\}\), \(B_3=\{3,4\}\), and so on.

For non-constant modes, a characteristic angular half-wavelength scale is
\[
  d_n(t)\sim \frac{\pi a(t)}{\mathcal F_n},
  \qquad n\ge1 .
\]
At the Planck boundary \(a_{\rm eff}(t_P)\simeq\ell_P\), this formula is a formal
spectral scale.  The effective theory does not claim operational access to
spatial resolution below \(\ell_P\).

\begin{proposition}[Planck-boundary resolution of the stratification]
\label{prop:resolution}
Assume \(a_{\rm eff}(t_P)\ge\ell_P\) and assume that the effective theory assigns
no operational support to sub-Planck spatial resolution.  Then the Fibonacci
bands remain a well-defined organization of the compact \(S^3\) spectrum, but
only bands associated with resolvable observables belong to the semiclassical
description.  Modes whose formal wavelength is below \(\ell_P\) are part of the
mathematical spectral tower, not independent resolved degrees of freedom of the
effective Planck-boundary theory.
\end{proposition}

\begin{remark}
The golden ratio enters through
\[
  \lim_{n\to\infty}\frac{\mathcal F_{n+1}}{\mathcal F_n}
  =\varphi=\frac{1+\sqrt5}{2}.
\]
This is a consequence of the Fibonacci recurrence.  It is not a new local
coupling, field, or physical constant.
\end{remark}

\subsection{Metric near the boundary}

The closed FLRW metric is
\[
  ds^2=-dt^2+a(t)^2
  \bigl[d\chi^2+\sin^2\chi\,(d\vartheta^2+\sin^2\vartheta\,d\varphi^2)\bigr].
\]
In the homogeneous semiclassical regime, \(a(t)\) is interpreted as an effective
scale variable, for instance as an expectation value of a geometric operator in
a suitable state.  Quantum corrections are represented schematically by
\[
  g_{ij}(t)=g^{(0)}_{ij}(t)+\hbar\,\Delta g_{ij}(t),
\]
where \(g^{(0)}_{ij}\) is the round \(S^3\) spatial metric with physical scale
factor \(a(t)\).  This notation records a semiclassical expansion; it is not a
microscopic quantisation prescription.

\begin{assumption}[Regular Planck-boundary semiclassical data]
\label{ass:regular-data}
At the Planck boundary the effective scale factor and kinematic data satisfy
\[
\begin{aligned}
  a_{\rm eff}(t_P)&\ge a_{\rm min}\ge \ell_P,\\
  |H(t_P)|&<\infty,\\
  |\dot H(t_P)|&<\infty .
\end{aligned}
\]
Metric fluctuations are assumed bounded in the restricted sense that their
contribution to \(a_{\rm eff}\), \(H\), and \(\dot H\) remains finite.
\end{assumption}

\begin{remark}
The spatial metric does not become Euclidean.  In the homogeneous \(k=+1\)
sector the background spatial geometry remains the round metric on \(S^3\),
regularised only by the non-zero effective radius.
\end{remark}

%% ================================================================
\section{Kinematic Curvature Bound}
%% ================================================================
\label{sec:curvature}

For a closed FLRW metric with \(k=+1\), the four-dimensional Kretschmann scalar is
\[
\begin{aligned}
  K
  &=R_{\mu\nu\rho\sigma}R^{\mu\nu\rho\sigma}  \\
  &=12\left[
      \left(\frac{\ddot a}{a}\right)^2
      +\left(\frac{\dot a^2+1}{a^2}\right)^2
    \right]                                      \\
  &=12\left[
      (\dot H+H^2)^2
      +\left(H^2+\frac1{a^2}\right)^2
    \right].
\end{aligned}
\]
Classically, \(a\to0\) together with uncontrolled \(H\) or \(\dot H\) produces a
curvature singularity.  The present framework instead removes \(a=0\) from the
effective configuration space.

\begin{theorem}[Kinematic curvature bound from regular Planck-boundary data]
\label{thm:curvature}
For a closed FLRW metric with \(k=+1\), assume that at the Planck boundary
\[
  a_{\rm eff}(t_P)=\alpha_P\ell_P,\qquad \alpha_P\ge1,
\]
and
\[
  |H(t_P)|<\infty,\qquad |\dot H(t_P)|<\infty .
\]
Then
\[
  K(t_P)\leq \frac{C}{\ell_P^4}<\infty,
\]
where
\[
\begin{aligned}
  C&=12\left[(q_P+h_P^2)^2+\left(h_P^2+\alpha_P^{-2}\right)^2\right],\\
  h_P&=\ell_P H(t_P),\qquad
  q_P=\ell_P^2\dot H(t_P).
\end{aligned}
\]
\end{theorem}

\begin{proof}
Substitution into the FLRW expression gives
\[
  K(t_P)=
  \frac{12}{\ell_P^4}
  \left[(q_P+h_P^2)^2+\left(h_P^2+\alpha_P^{-2}\right)^2\right].
\]
The assumptions imply that \(\alpha_P^{-1}\), \(h_P\), and \(q_P\) are finite.
Therefore \(C\) is finite.
\end{proof}

\begin{remark}
The theorem is not a dynamical singularity-resolution theorem.  It states the
exact kinematic consequence of regular boundary data.  A complete quantum-gravity
model would still have to explain why the finite values of \(H(t_P)\) and
\(\dot H(t_P)\) are dynamically selected.
\end{remark}

\begin{corollary}
Within any interval on which \(a_{\rm eff}\ge\ell_P\) and \(H,\dot H\) remain
finite, the closed-FLRW curvature scalar \(K\) remains finite.
\end{corollary}

\begin{interpretation}[No external Planck-boundary time]
\label{int:no-external-time}
The Wheeler--DeWitt equation is a constraint equation, \(\hat H\Psi=0\), and
does not contain an external time parameter.  A semiclassical clock can emerge in
a WKB/Born--Oppenheimer regime, but that approximation is not assumed at the
Planck boundary.  Thus statements about \(\partial_t g_{ij}\) at
\(a\simeq\ell_P\) are clock- and gauge-dependent unless a specific clock
reconstruction has been supplied.
\end{interpretation}

%% ================================================================
\section{Spectral Zeta Regularisation of the Vacuum Scale}
%% ================================================================

\subsection{Scale-covariant compact spectra}

For a free field on a compact spatial manifold the formal vacuum energy is
\[
  E_{\rm vac}=\frac12\sum_k\omega_k .
\]
On a round \(S^3\) of physical radius \(a\), a scale-covariant spectral family
has
\[
  \omega_k(a)=\frac{\nu_k}{a},
\]
where \(\nu_k\) is the dimensionless spectrum on the unit sphere.  Define
\[
  \zeta_\omega(s)=\sum_k\nu_k^{-s},
  \qquad
  \zeta_\omega(s;a)=\sum_k\omega_k(a)^{-s}=a^s\zeta_\omega(s).
\]
If the spectral zeta function admits a finite analytic continuation to \(s=-1\),
the zeta-regularised Casimir energy is
\[
  E_{\rm vac}^{\rm reg}(a)=\frac{1}{2a}\zeta_\omega(-1)
  =\frac{\Czeta}{a},
  \qquad
  \Czeta=\frac12\zeta_\omega(-1).
\]
Since
\[
  V_{S^3}(a)=2\pi^2a^3,
\]
the corresponding regularised density is
\[
  \rho_{\rm vac}^{\rm reg}(a)
  =\frac{\Czeta}{2\pi^2a^4}.
\]
The \(a^{-4}\) law follows from scale covariance: energy scales as \(a^{-1}\)
and volume as \(a^3\).

\subsection{Conformally shifted scalar reference spectrum}

For a massless conformally shifted scalar on \(S^3\), the relevant unit-sphere
operator is \(-\Delta_{S^3}+1\), with
\[
  \nu_\ell=\ell+1,\qquad d_\ell=(\ell+1)^2,\qquad
  \ell=0,1,2,\ldots .
\]
Thus
\[
  \zeta_\omega(s)
  =\sum_{\ell=0}^\infty(\ell+1)^2(\ell+1)^{-s}
  =\sum_{n=1}^\infty n^{2-s}
  =\Riemann(s-2).
\]
Analytic continuation gives
\[
  \zeta_\omega(-1)=\Riemann(-3)=\frac1{120},
  \qquad
  \Czeta=\frac1{240},
\]
and therefore
\[
  \rho_{\rm vac}^{\rm reg}(a)=\frac{1}{480\pi^2a^4}.
\]

\begin{remark}
The numerical coefficient and sign of a Casimir contribution depend on field
content, spin structure, operator choice, subtraction prescription, and
renormalisation convention.  In semiclassical gravity the absolute vacuum energy
also mixes with the renormalisation of the cosmological and gravitational
couplings.  The invariant statement used here is the compact spectral scaling
and finiteness for \(a\ge\ell_P\), not a direct prediction of the observed
cosmological constant.
\end{remark}

\subsection{Fibonacci organization and the zeta sum}

The Fibonacci bands \(B_n\) organize the same \(S^3\) spectral tower.  For a
characteristic dimensionless frequency in a band one may use, for example,
\[
  \Omega_n=\sqrt{\nu_{\mathcal F_n}\nu_{\mathcal F_{n+1}-1}},
  \qquad n\ge1 .
\]
Then
\[
  \frac{\Omega_{n+1}}{\Omega_n}\longrightarrow\varphi
  \qquad (n\to\infty),
\]
provided the underlying spectrum is asymptotically linear in \(\ell\), as it is
for the conformal scalar frequencies.

A schematic band-counting zeta
\[
  \zeta_{\rm band}(s)=\sum_{n\ge1}|B_n|\Omega_n^{-s}
\]
behaves asymptotically like a geometric series in \(\varphi^{1-s}\), so it
converges for \(\operatorname{Re}s>1\).  If the full degeneracy of the
three-dimensional spectrum is included, the large-band contribution scales as
\(\mathcal F_n^{3-s}\), and the usual convergence domain
\(\operatorname{Re}s>3\) is recovered before analytic continuation.

\begin{figure}[!tbp]
\centering
\includegraphics[width=0.98\columnwidth]{figures/fibonacci_ratio_convergence.pdf}
\caption{The ratio \(\mathcal F_{n+1}/\mathcal F_n\) converges to
\(\varphi\).  The golden ratio is an asymptotic coding scale of the Fibonacci
coarse-graining, not an independent physical coupling.}
\label{fig:fibonacci-ratio}
\end{figure}

\begin{theorem}[Spectral-zeta vacuum-density scale]
\label{thm:vacenergy}
Let a compact \(S^3\) spectral family satisfy \(\omega_k(a)=\nu_k/a\), and let
\(\zeta_\omega(s)=\sum_k\nu_k^{-s}\) admit a finite analytic continuation at
\(s=-1\).  Then
\[
\begin{aligned}
  E_{\rm vac}^{\rm reg}(a)&=\frac{\Czeta}{a},\\
  \rho_{\rm vac}^{\rm reg}(a)&=\frac{\Czeta}{2\pi^2a^4},\\
  \Czeta&=\frac12\zeta_\omega(-1).
\end{aligned}
\]
In particular, \(\rho_{\rm vac}^{\rm reg}\) is finite for every
\(a\ge\ell_P\).  For the conformally shifted scalar reference spectrum,
\[
  \rho_{\rm vac}^{\rm reg}(a)=\frac{1}{480\pi^2a^4}.
\]
\end{theorem}

\begin{proof}
Scale covariance gives
\[
  E_{\rm vac}^{\rm reg}(a)
  =\frac12\sum_k^{\rm reg}\frac{\nu_k}{a}
  =\frac{1}{2a}\zeta_\omega(-1)
  =\frac{\Czeta}{a}.
\]
Dividing by \(V_{S^3}(a)=2\pi^2a^3\) gives the stated density.  Since
\(a\ge\ell_P>0\) and \(\zeta_\omega(-1)\) is finite by hypothesis, the density is
finite at the Planck boundary.
\end{proof}

\begin{figure}[!tbp]
\centering
\includegraphics[width=0.98\columnwidth]{figures/vacuum_density_scaling.pdf}
\caption{Dimensionless compact-spectrum vacuum-density scale for the conformally
shifted scalar reference spectrum.  The exact scaling is \(a^{-4}\).}
\label{fig:vacuum-scaling}
\end{figure}

%% ================================================================
\section{Self-Adjoint Wheeler--DeWitt Boundary}
%% ================================================================

The minisuperspace wave function is defined on the effective domain
\[
  a\in[\ell_P,\infty).
\]
The sub-Planck interval \(0<a<\ell_P\) is not part of the Hilbert-space measure
of this effective theory.

A schematic factor ordering for the Wheeler--DeWitt operator is
\[
  \hat H
  =-\frac{\hbar^2}{12\pi\mP^2}\frac1a
    \frac{\partial}{\partial a}
    \left(a\frac{\partial}{\partial a}\right)
    +\hat H_m+U_{\rm eff}(a),
\]
where \(U_{\rm eff}\) contains curvature, matter, and regularised compact
spectral contributions.  The later argument requires not this precise potential
but self-adjointness of the operator on the chosen domain.

The minisuperspace Hilbert space is
\[
\begin{aligned}
  \mathcal H_{\rm mini}&=L^2([\ell_P,\infty),a\,da),\\
  \langle\Phi|\Psi\rangle
  &=\int_{\ell_P}^{\infty}a\,\Phi^*(a)\Psi(a)\,da .
\end{aligned}
\]
At the regular endpoint \(a=\ell_P\), a local one-parameter self-adjoint boundary
condition is
\[
  \cos\alpha\,\Psi(\ell_P)
  +\sin\alpha\,\ell_P\Psi'(\ell_P)=0,
  \qquad \alpha\in[0,\pi).
\]
At \(a=\infty\) we assume either the usual limit-point case or, more generally,
that the operator domain is chosen so that the boundary form vanishes there.
Thus no probability flux is allowed through either endpoint of the effective
half-line.

\begin{theorem}[Norm preservation in the effective minisuperspace domain]
\label{thm:unitarity}
Assume that the Wheeler--DeWitt operator is given a self-adjoint extension on
\(L^2([\ell_P,\infty),a\,da)\), with vanishing boundary form at
\(a=\ell_P\) and at \(a=\infty\).  Then, for every relational or semiclassical
clock parameter \(\tau\) for which an effective Schr\"odinger form is valid,
\[
  \frac{d}{d\tau}\|\Psi\|^2=0 .
\]
Equivalently, no probability is lost into the excluded sub-Planck region.
\end{theorem}

\begin{proof}
For the Sturm--Liouville differential expression
\[
  L\Psi=-\frac1a(a\Psi')',
\]
the boundary form is proportional to
\[
  \bigl[a(\Phi^*\Psi'-\Phi'^*\Psi)\bigr]_{\ell_P}^{\infty}.
\]
Self-adjointness is precisely the vanishing of this form on the operator domain.
Taking \(\Phi=\Psi\) gives zero probability flux at the endpoints.  The usual
Schr\"odinger conservation identity in any valid relational-time reconstruction
therefore yields \(d\|\Psi\|^2/d\tau=0\).
\end{proof}

\begin{corollary}[Regular Planck boundary]
The boundary condition permits \(\Psi(\ell_P)\ne0\) for generic Robin parameter
\(\alpha\).  The Planck boundary is therefore a regular endpoint of the
effective quantum problem, not an absorbing singularity.
\end{corollary}

\subsection{WKB behaviour near the boundary}

If a minisuperspace potential produces a classically forbidden interval adjacent
to the Planck boundary, a WKB branch may have the schematic form
\[
  \Psi(a)\sim A\exp\!\left[
  -\frac1\hbar\int_{\ell_P}^{a}|p_a(a')|\,da'
  \right],
  \qquad a\ge\ell_P .
\]
This expression is used only on the physical domain.  The value at
\(a=\ell_P\) is determined by the self-adjoint boundary condition and need not
vanish.

\begin{figure}[!tbp]
\centering
\includegraphics[width=0.98\columnwidth]{figures/planck_boundary_wavefunction.pdf}
\caption{Schematic finite-amplitude WKB profile on the effective domain
\(a\ge\ell_P\).  No probability measure is assigned to \(0<a<\ell_P\).}
\label{fig:planck-wavefunction}
\end{figure}

\begin{corollary}[State-dependent Planck-scale correlations]
In a compact \(S^3\) universe, globally constrained quantum states need not
factorise over arbitrary spatial partitions.  Correlations between subregions
may therefore persist to the Planck boundary.  The corresponding entanglement
entropy is state-dependent and is not fixed by topology alone.
\end{corollary}

%% ================================================================
\section{Noether-II Identities on Closed Cauchy Slices}
%% ================================================================

Let the generally covariant action be written schematically as
\[
  S[g,\psi]=\int_{\mathcal M}d^4x\,\sqrt{-g}\,
  \mathcal L(g_{\mu\nu},\psi,\nabla\psi).
\]
Diffeomorphism invariance is a local gauge symmetry.  By Noether's second
theorem~\cite{Noether}, it gives differential identities among the Euler--Lagrange
expressions.  It does not give, by itself, an independent global energy
conservation law of the Noether-I type.

\begin{theorem}[Noether-II identities and closed-slice boundary charges]
\label{thm:noether}
Let a generally covariant theory be defined on a spacetime with closed Cauchy
slice \(\Sigma\simeq S^3\).  Then:
\begin{enumerate}
\item Diffeomorphism invariance gives Noether-II differential identities among
the Euler--Lagrange expressions.  In Einstein gravity coupled to matter these
include the contracted Bianchi identity and the corresponding covariant
stress-energy conservation law on shell.
\item In Iyer--Wald covariant phase-space notation~\cite{IyerWald},
\[
  J_\xi=\Theta(\phi;\mathcal L_\xi\phi)-\xi\cdot L,
  \qquad
  J_\xi\simeq dQ_\xi ,
\]
and the Hamiltonian variation associated with \(\xi\) is
\[
  \delta H_\xi
  =\int_{\partial\Sigma}
  \left(\delta Q_\xi-\xi\cdot\Theta\right).
\]
For the full closed slice, \(\partial\Sigma=\varnothing\), so
\[
  \delta H_\xi=0
\]
for pure-gauge diffeomorphisms on the full slice.
\item This full-slice statement does not exclude quasi-local charges on internal
subdomains \(D\subset\Sigma\) with \(\partial D\ne\varnothing\).
\end{enumerate}
\end{theorem}

\begin{proof}
Part (1) is the standard Noether-II consequence of local diffeomorphism
invariance.  Part (2) is the Iyer--Wald variational identity evaluated on a
Cauchy slice with empty boundary.  Part (3) follows because an internal subdomain
can have a boundary even when the whole Cauchy slice does not.
\end{proof}

\begin{remark}
For compact cosmology there is no ADM boundary at spatial infinity.  Global
``energy'' is therefore not a boundary charge of the asymptotically flat type.
The Hamiltonian content of general relativity in the closed case is carried by
constraints and possible internal or quasi-local boundary terms.
\end{remark}

\subsection{Asymptotic Fibonacci scaling}

The Fibonacci shift \(n\mapsto n+1\) is discrete and only asymptotically
represented by multiplication by \(\varphi\).  It is therefore not a continuous
variational symmetry and does not fall under Noether's first or second theorem.

\begin{proposition}[Asymptotic Fibonacci scaling law]
\label{prop:fib-scaling}
For an exactly scale-covariant spectral family,
\[
  \zeta_\omega(s;a)=a^s\zeta_\omega(s).
\]
Consequently,
\[
  \zeta_\omega(s;\varphi a)=\varphi^s\zeta_\omega(s;a),
\]
and at the regularised Casimir point,
\[
  E_{\rm vac}^{\rm reg}(\varphi a)
  =\varphi^{-1}E_{\rm vac}^{\rm reg}(a),
  \qquad
  \rho_{\rm vac}^{\rm reg}(\varphi a)
  =\varphi^{-4}\rho_{\rm vac}^{\rm reg}(a).
\]
The Fibonacci interpretation of this rescaling is asymptotic in the band index.
\end{proposition}

\begin{proof}
The first identity is the definition of scale covariance.  Substituting
\(a\mapsto\varphi a\) gives the second.  The energy scaling follows at
\(s=-1\), and the density acquires the additional volume factor
\((\varphi a)^{-3}\).
\end{proof}

%% ================================================================
\section{Possible Observational Signatures}
%% ================================================================

The following items are qualitative targets, not predictions derived in this
paper.  A connection between Planck-boundary data and observations requires a
separate perturbation-transfer model.

\medskip
\noindent\textbf{CMB power spectrum.}
The Fibonacci banding motivates searches for residual structure near multipoles
\[
  \ell=2,3,5,8,13,21,34,55,89,\ldots .
\]
The dipole \(\ell=1\) is not a clean primordial observable because it is
dominated by the observer's motion.  No universal amplitude is derived here.

\medskip
\noindent\textbf{Primordial power spectrum.}
A regular Planck boundary can modify admissible initial data for perturbation
modes.  The scale \(1/\ell_P\) is ultraviolet, while low-\(\ell\) CMB anomalies
are infrared; relating them requires explicit inflationary or pre-inflationary
transfer dynamics.

\medskip
\noindent\textbf{Dark-energy sector.}
The compact-spectrum term obtained above scales as \(a^{-4}\).  It therefore
does not by itself mimic a constant late-time cosmological constant.  DESI-class
evolving-dark-energy fits are external phenomenological targets, not predictions
of the present theorem~\cite{DESI}.

\medskip
\noindent\textbf{Stochastic gravitational-wave background.}
A discrete compact \(S^3\) tensor spectrum could motivate searches for resonance
or modulation features in a stochastic background.  The present-day frequencies
would depend on the full expansion and transfer history and are not fixed here.

\begin{table*}[!tbp]
\centering
\small
\renewcommand{\arraystretch}{1.25}
\caption{Possible qualitative signatures.  None is a definitive prediction of
the present Planck-boundary theorem.}
\label{tab:signatures}
\begin{tabular}{p{0.25\textwidth}p{0.43\textwidth}p{0.22\textwidth}}
\toprule
Observable & Possible qualitative signature & Required extra input\\
\midrule
CMB low multipoles & Residual structure near Fibonacci multipoles \(\ell\ge2\) &
Perturbation transfer model\\
Primordial spectrum & Modified admissible initial data & Inflationary or pre-inflationary dynamics\\
Dark-energy sector & No direct prediction from the \(a^{-4}\) term & Additional late-time dynamics\\
Compact Casimir component & \(\rho_{\rm vac}^{\rm reg}\propto a^{-4}\) &
Renormalisation and field content\\
SGWB & Possible discrete-spectrum modulation & Expansion and tensor-mode transfer history\\
\bottomrule
\end{tabular}
\end{table*}

%% ================================================================
\section{Discussion}
%% ================================================================

The model proves conditional statements inside a deliberately narrow domain.  Its
logical structure is
\[
\begin{aligned}
  &\text{closed }S^3\text{ sector}\\
  &+ a\ge \ell_P\\
  &+ \text{finite }H,\dot H\\
  &+ \text{self-adjoint WDW boundary}\\
  &+ \text{compact spectral zeta regularisation}.
\end{aligned}
\]
From these inputs one obtains finite FLRW curvature at the boundary, finite
regularised compact-spectrum vacuum density, norm preservation in the
minisuperspace Hilbert space, and the standard Noether-II identities of
generally covariant dynamics.

\subsection{Claims established here}

\begin{enumerate}
\item Within the smooth closed connected orientable simply-connected spatial
sector, the Poincar\'e--Perelman theorem together with Moise's theorem selects
the standard \(S^3\) sector.
\item Under finite Planck-boundary data, the closed-FLRW Kretschmann scalar
obeys \(K(t_P)\le C/\ell_P^4\).
\item Zeta regularisation of a scale-covariant compact \(S^3\) spectrum gives
\(\rho_{\rm vac}^{\rm reg}=\Czeta/(2\pi^2a^4)\), finite for
\(a\ge\ell_P\).
\item A self-adjoint Wheeler--DeWitt boundary condition prevents probability
flux through \(a=\ell_P\) in the effective minisuperspace description.
\item Diffeomorphism invariance gives Noether-II differential identities, and
the full closed Cauchy slice has no Iyer--Wald boundary Hamiltonian variation.
\end{enumerate}

\subsection{Claims not established here}

The following are not consequences of the present model:
\begin{enumerate}
\item a microscopic derivation of the Planck boundary;
\item a derivation of finite \(H(t_P)\) and \(\dot H(t_P)\);
\item an LQC-type bounce;
\item a prediction of the observed cosmological constant;
\item an exact Fibonacci spacetime symmetry;
\item a universal entanglement entropy fixed by topology;
\item a quantitative CMB or SGWB prediction.
\end{enumerate}

\subsection{Relation to loop quantum cosmology and CDT}

The curvature result resembles loop quantum cosmology in its final boundedness
statement, but the mechanism is different.  LQC derives effective boundedness
from polymer quantisation and a discrete geometric spectrum~\cite{Ashtekar}.
Here boundedness is conditional on a self-adjoint Planck boundary and finite
semiclassical boundary data.

Causal dynamical triangulations supplies another independent example in which
Planck-scale geometry may differ from smooth classical intuition~\cite{CDT}.
The present Fibonacci coarse-graining is only an analogy at the level of
effective spectral organization; it is not a derivation of CDT spectral
dimensional flow.

%% ================================================================
\section{Summary}
%% ================================================================

The \model{} framework gives a rigorous conditional Planck-boundary construction
for compact \(S^3\) minisuperspace cosmology.  Its main results are:

\begin{enumerate}
\item \textbf{Topological sector.}
Assuming smooth closed connected orientable simple connectivity, the spatial
slice is the standard \(S^3\).

\item \textbf{Fibonacci coarse-graining.}
The \(S^3\) Laplacian spectrum may be organized into Fibonacci bands.  The
golden ratio enters only through the asymptotic limit
\(\mathcal F_{n+1}/\mathcal F_n\to\varphi\).

\item \textbf{Curvature bound.}
For closed FLRW data with \(a_{\rm eff}(t_P)=\alpha_P\ell_P\), \(\alpha_P\ge1\),
and finite \(H(t_P)\), \(\dot H(t_P)\),
\[
  K(t_P)\le\frac{C}{\ell_P^4}<\infty .
\]

\item \textbf{Vacuum scale.}
For a scale-covariant compact spectrum with finite \(\zeta_\omega(-1)\),
\[
  \rho_{\rm vac}^{\rm reg}(a)=\frac{\Czeta}{2\pi^2a^4}.
\]
For the conformally shifted scalar,
\[
  \rho_{\rm vac}^{\rm reg}(a)=\frac{1}{480\pi^2a^4}.
\]

\item \textbf{Self-adjoint boundary.}
A self-adjoint Wheeler--DeWitt operator on
\(L^2([\ell_P,\infty),a\,da)\) gives zero probability flux through the Planck
boundary.

\item \textbf{Noether-II structure.}
Diffeomorphism invariance gives differential identities.  On the full closed
slice, Iyer--Wald boundary Hamiltonian variations vanish because
\(\partial\Sigma=\varnothing\), while internal quasi-local charges are not
excluded.
\end{enumerate}

Thus the model is mathematically coherent as a conditional effective
Planck-boundary framework.  Its strength is not that it solves quantum gravity,
but that it states exactly which assumptions are sufficient for bounded
curvature, compact-spectrum finiteness, and minisuperspace norm preservation.

%% ================================================================
\begin{thebibliography}{99}
%% ================================================================

\bibitem{Perelman1}
G.~Perelman,
The entropy formula for the Ricci flow and its geometric applications,
arXiv:math/0211159 (2002).

\bibitem{Perelman2}
G.~Perelman,
Ricci flow with surgery on three-manifolds,
arXiv:math/0303109 (2003).

\bibitem{Perelman3}
G.~Perelman,
Finite extinction time for the solutions to the Ricci flow on certain three-manifolds,
arXiv:math/0307245 (2003).

\bibitem{Hamilton}
R.~S.~Hamilton,
Three-manifolds with positive Ricci curvature,
\textit{J. Differential Geom.} \textbf{17}, 255 (1982).

\bibitem{Moise}
E.~E.~Moise,
Affine structures in 3-manifolds. V. The triangulation theorem and Hauptvermutung,
\textit{Ann. of Math.} \textbf{56}, 96 (1952).

\bibitem{Zeckendorf}
E.~Zeckendorf,
Repr\'esentation des nombres naturels par une somme de nombres de Fibonacci,
\textit{Bull. Soc. R. Sci. Li\`ege} \textbf{41}, 179 (1972).

\bibitem{Elizalde}
E.~Elizalde, S.~D.~Odintsov, A.~Romeo, A.~A.~Bytsenko, and S.~Zerbini,
\textit{Zeta Regularization Techniques with Applications}
(World Scientific, Singapore, 1994).

\bibitem{Vassilevich}
D.~V.~Vassilevich,
Heat kernel expansion: User's manual,
\textit{Phys. Rep.} \textbf{388}, 279 (2003),
arXiv:hep-th/0306138.

\bibitem{DeWitt}
B.~S.~DeWitt,
Quantum theory of gravity. I. The canonical theory,
\textit{Phys. Rev.} \textbf{160}, 1113 (1967).

\bibitem{Wheeler}
J.~A.~Wheeler,
Superspace and the nature of quantum geometrodynamics,
in \textit{Batelle Rencontres}, Benjamin, New York, pp.~242--307 (1968).

\bibitem{Hartle}
J.~B.~Hartle and S.~W.~Hawking,
Wave function of the Universe,
\textit{Phys. Rev. D} \textbf{28}, 2960 (1983).

\bibitem{PageWootters}
D.~N.~Page and W.~K.~Wootters,
Evolution without evolution,
\textit{Phys. Rev. D} \textbf{27}, 2885 (1983).

\bibitem{Isham}
C.~J.~Isham,
Canonical quantum gravity and the problem of time,
arXiv:gr-qc/9210011 (1993).

\bibitem{Noether}
E.~Noether,
Invariante Variationsprobleme,
\textit{Nachr. Ges. Wiss. G\"ottingen} \textbf{1918}, 235 (1918).

\bibitem{Wald}
R.~M.~Wald,
\textit{General Relativity}
(University of Chicago Press, Chicago, 1984).

\bibitem{IyerWald}
V.~Iyer and R.~M.~Wald,
Some properties of Noether charge and a proposal for dynamical black hole entropy,
\textit{Phys. Rev. D} \textbf{50}, 846 (1994),
arXiv:gr-qc/9403028.

\bibitem{Ashtekar}
A.~Ashtekar, T.~Pawlowski, and P.~Singh,
Quantum nature of the Big Bang,
\textit{Phys. Rev. Lett.} \textbf{96}, 141301 (2006),
arXiv:gr-qc/0602086.

\bibitem{CDT}
J.~Ambj\o{}rn, J.~Jurkiewicz, and R.~Loll,
Emergence of a 4D world from causal quantum gravity,
\textit{Phys. Rev. Lett.} \textbf{95}, 171301 (2005),
arXiv:hep-th/0505113.

\bibitem{Planck2018}
Planck Collaboration,
Planck 2018 results. VI. Cosmological parameters,
\textit{Astron. Astrophys.} \textbf{641}, A6 (2020),
arXiv:1807.06209.

\bibitem{DESI}
M.~Abdul Karim et al. (DESI Collaboration),
DESI DR2 results. II. Measurements of baryon acoustic oscillations and cosmological constraints,
\textit{Phys. Rev. D} \textbf{112}, 083515 (2025),
arXiv:2503.14738, doi:10.1103/tr6y-kpc6.

\bibitem{Carroll}
S.~M.~Carroll,
\textit{Spacetime and Geometry: An Introduction to General Relativity}
(Addison-Wesley, San Francisco, 2004).

\bibitem{Weinberg}
S.~Weinberg,
\textit{Cosmology}
(Oxford University Press, Oxford, 2008).

\end{thebibliography}

\end{document}
